\section{Implement Upper Bound}
\label{sec:bs-upper-bound}

\problemstatement{We are given a sorted array $\textit{arr}$ of $n$
elements and a value $X$. We must find the \textbf{upper bound} of $X$:
the index of the first element in $\textit{arr}$ that is
\textbf{strictly greater than} $X$. If no such element exists, the upper
bound is $n$.}

\begin{patternbox}
Upper bound is lower bound's twin, with one word changed: ``greater than or
equal to'' becomes ``strictly greater than.'' Whenever you see \emph{``first
index strictly greater than $X$''} as opposed to \emph{``first index at
least $X$,''} you are being asked for an upper bound rather than a lower
bound. The two are easy to confuse under time pressure, so it is worth
fixing the distinction precisely: lower bound finds where $X$ \emph{could
first appear}; upper bound finds where the elements equal to $X$
\emph{end}.
\end{patternbox}

\subsection*{Understanding the problem carefully}

The predicate here is $P(i) = (\textit{arr}[i] > X)$, which, exactly as in
lower bound, is \texttt{false} for a prefix and \texttt{true} for a suffix
of the sorted array, with a single crossover. The only change from
Section~\ref{sec:bs-lower-bound} is the strictness of the comparison: an
element exactly equal to $X$ satisfies the lower-bound predicate
($\ge X$) but \emph{not} the upper-bound predicate ($> X$). Consequently,
if $X$ appears multiple times in the array, the lower bound points to the
first occurrence of $X$, while the upper bound points to the index
\emph{immediately after} the last occurrence of $X$ -- a fact we will use
directly in Section~\ref{sec:bs-first-last} and
Section~\ref{sec:bs-count-occurrences}.

\begin{ideabox}[Candidate Space and Predicate]
The candidate space is indices $0, \dots, n$. The predicate is
$P(i) = (\textit{arr}[i] > X)$ for $i < n$, and we want the smallest index
where $P$ holds.
\end{ideabox}

\subsection*{Why the same template applies, with one flipped comparison}

The reasoning is identical to lower bound: whenever
$\textit{arr}[\textit{mid}] > X$, this index is a valid candidate for the
upper bound, but there might be an even earlier index that is also strictly
greater than $X$, so we record it and search left. Whenever
$\textit{arr}[\textit{mid}] \le X$ (note: this now includes equality, unlike
lower bound), this index and everything before it is disqualified, so we
search right. The only code change relative to lower bound is replacing the
comparison \texttt{arr[mid] >= X} with \texttt{arr[mid] > X}.

\subsection*{Algorithm}

\begin{enumerate}
  \item Initialize $\textit{lo} \gets 0$, $\textit{hi} \gets n-1$,
        $\textit{ans} \gets n$.
  \item While $\textit{lo} \le \textit{hi}$:
  \begin{enumerate}
    \item $\textit{mid} \gets \textit{lo} + (\textit{hi}-\textit{lo})/2$.
    \item If $\textit{arr}[\textit{mid}] > X$: record
          $\textit{ans} \gets \textit{mid}$; $\textit{hi} \gets
          \textit{mid}-1$.
    \item Otherwise: $\textit{lo} \gets \textit{mid}+1$.
  \end{enumerate}
  \item Return \textit{ans}.
\end{enumerate}

\begin{algorithm}[H]
\caption{Upper Bound}
\begin{algorithmic}[1]
\Procedure{UpperBound}{$\textit{arr}, X$}
  \State $\textit{lo} \gets 0$, $\textit{hi} \gets n-1$, $\textit{ans} \gets n$
  \While{$\textit{lo} \le \textit{hi}$}
    \State $\textit{mid} \gets \textit{lo} + (\textit{hi}-\textit{lo})/2$
    \If{$\textit{arr}[\textit{mid}] > X$}
      \State $\textit{ans} \gets \textit{mid}$
      \State $\textit{hi} \gets \textit{mid} - 1$
    \Else
      \State $\textit{lo} \gets \textit{mid} + 1$
    \EndIf
  \EndWhile
  \State \Return \textit{ans}
\EndProcedure
\end{algorithmic}
\end{algorithm}

\subsection*{Dry run}

\dryrun{Take $\textit{arr} = [1, 2, 2, 3]$, $X = 2$.}

\begin{itemize}
  \item $\textit{lo}=0,\textit{hi}=3,\textit{ans}=4$: $\textit{mid}=1$,
        $\textit{arr}[1]=2 > 2$? No. $\textit{lo}=2$.
  \item $\textit{lo}=2,\textit{hi}=3$: $\textit{mid}=2$,
        $\textit{arr}[2]=2 > 2$? No. $\textit{lo}=3$.
  \item $\textit{lo}=3,\textit{hi}=3$: $\textit{mid}=3$,
        $\textit{arr}[3]=3 > 2$? Yes. $\textit{ans}=3$, $\textit{hi}=2$.
  \item $\textit{lo}=3 > \textit{hi}=2$: loop ends.
\end{itemize}

The upper bound is index $3$ -- one past the last occurrence of $2$
(which sits at index $2$), consistent with the lower bound of $2$ found in
Section~\ref{sec:bs-lower-bound} being index $1$: the occurrences of $2$
span indices $[1, 2]$.

\subsection*{C++ implementation}

\begin{lstlisting}
#include <bits/stdc++.h>
using namespace std;

int upperBound(vector<int>& arr, int X) {
    int n = arr.size();
    int lo = 0, hi = n - 1, ans = n;

    while (lo <= hi) {
        int mid = lo + (hi - lo) / 2;
        if (arr[mid] > X) {
            ans = mid;
            hi = mid - 1;
        } else {
            lo = mid + 1;
        }
    }
    return ans;
}

int main() {
    vector<int> arr = {1, 2, 2, 3};
    cout << "Upper bound of 2: " << upperBound(arr, 2) << endl;
    return 0;
}
\end{lstlisting}

\begin{complexitybox}
\textbf{Time Complexity.} $O(\log n)$, identical to lower bound, since
the algorithm structure is unchanged.
\textbf{Space Complexity.} $O(1)$ auxiliary space.
\end{complexitybox}

\begin{takeawaybox}
Lower bound and upper bound differ by exactly one comparison operator
($\ge$ versus $>$), but that single difference is what makes the pair
together so useful: the number of occurrences of $X$ in a sorted array is
always $\textit{upperBound}(X) - \textit{lowerBound}(X)$, and the first and
last occurrence indices of $X$ are $\textit{lowerBound}(X)$ and
$\textit{upperBound}(X) - 1$ respectively (when $X$ is present at all).
\end{takeawaybox}
